\documentclass[10pt]{article}

\usepackage[margin=0.75in]{geometry}
\usepackage{amsmath,amsthm,amssymb,bm}
\usepackage{xcolor}
\usepackage{cancel}
\usepackage{graphicx}
\usepackage{changepage}
\usepackage{circuitikz}
\usepackage{pgfplots}
\usepackage{physics}
\usepackage{hyperref}
\usepackage{siunitx}
\usepackage{minted}
\usepackage{fontspec}
\usepackage{relsize}
\usepackage{subfig}
\usepackage{mhchem}
\usepackage{todonotes}
\usepackage{biblatex}
\usepackage{multicol, multirow, booktabs}
\usepackage[breakable]{tcolorbox}
\usepackage[inline]{enumitem}
\patchcmd{\thebibliography}{\section*{\refname}}{}{}{}
\addbibresource{sources.bib}

\theoremstyle{definition}
\newtheorem{problem}{Problem}
\newtheorem{soln}{Solution}

\pgfplotsset{compat=newest}
\usetikzlibrary{lindenmayersystems}
\usetikzlibrary{arrows}
\usetikzlibrary{calc}
\usetikzlibrary{positioning, fit}
\usetikzlibrary{3d, perspective}

\definecolor{incolor}{HTML}{303F9F}
\definecolor{outcolor}{HTML}{D84315}
\definecolor{cellborder}{HTML}{CFCFCF}
\definecolor{cellbackground}{HTML}{F7F7F7}
\newcommand{\ui}{\hat{i}}
\newcommand{\uj}{\hat{j}}
\newcommand{\uk}{\hat{k}}
\newcommand{\ux}{\hat{x}}
\newcommand{\uy}{\hat{y}}
\newcommand{\uz}{\hat{z}}
\newcommand{\uv}[1]{\hat{\bm{{#1}}}}
\newcommand{\pr}[1]{{#1^\prime}}
\newcommand{\justif}[2]{&{#1}&\text{#2}}
\newcommand{\dul}[1]{\underline{\underline{#1}}}
\newcommand{\pdiff}[2][]{{\frac{\partial^{#1}}{\partial {{#2}^{#1}}}}}
% \newcommand{\dif}[2][]{{\frac{d^{#1}}{d{{#2}^{#1}}}}}
\pgfdeclarelayer{background}  
\pgfsetlayers{background,main}
\AtBeginDocument{\RenewCommandCopy\qty\SI}

\makeatletter
\newcommand{\boxspacing}{\kern\kvtcb@left@rule\kern\kvtcb@boxsep}
\makeatother
\newcommand{\prompt}[4]{
    \ttfamily\llap{{\color{#2}[#3]:\hspace{3pt}#4}}\vspace{-\baselineskip}
}

\newcommand{\thevenin}[2]{
  \begin{center}
    \begin{circuitikz} \draw
      (0,0) -- (2,0) to[battery1, l_=$V_{Th}\eq#1$] (2,2) 
      to[resistor, l_=$R_{Th}\eq#2$] (0,2)
      ;
      \draw [o-] (-.07,2.079);
      \draw [o-] (-.07,0.079);
    \end{circuitikz}
  \end{center}
}

\newcommand{\norton}[2]{
  \begin{center}
    \begin{circuitikz} \draw
      (0,0) -- (3,0) to[american current source, l_=$I_{N}\eq#1$] (3,2) -- (0,2) (2,0)
      to[resistor, l=$R_{N}\eq#2$] (2,2)
      ;
      \draw [o-] (-.07,2.079);
      \draw [o-] (-.07,0.079);
    \end{circuitikz}
  \end{center}
}

\newcommand{\highlight}[1]{\colorbox{yellow}{#1}}

\newcommand{\ti}[1]{\widetilde{#1}}

\newfontface{\Kaufmann}{Kaufmann}
\DeclareTextFontCommand{\kf}{\Kaufmann}
\newcommand{\scriptr}{\fontsize{12pt}{12pt}\kf{r}}

\newfontface{\KaufmannB}{Kaufmann Bold}
\DeclareTextFontCommand{\kfb}{\KaufmannB}
\newcommand{\bscriptr}{\fontsize{12pt}{12pt}\kfb{r}}

\newcommand{\bv}[1]{\bm{\mathrm{#1}}}

\title{Physics 4605H: Project Outline}
\author{Jeremy Favro (0805980) \\ Trent University, Peterborough, ON, Canada}
\date{\today}

\begin{document}
\maketitle
My article choice is available at \url{https://doi.org/10.1103/PhysRevLett.109.050505}.


My overall goal is to focus on the fit error estimation because it's a uniquely (as far as I can tell) quantum computer problem. To do this I will probably need
to talk about the actual fitting algorithm, though I will have to skip much of the detail if I do. I think the way I would do this is to
try to construct a flow chart of the steps with some kind of very small scale/pictorial representation of what each step is doing. I'm not sure if these representations
would be built to illustrate the process with more of a focus on the quantum computing part or the algorithm.

Once I've illustrated how the fitting algorithm (FA) works I'll probably also have defined a lot of the terms I'll need
to talk about the error estimation algorithm (EEA), i.e. \emph{sparse matrix}, \emph{spectral norm}, \emph{Moore-Penrose Pseudoinverse}. From here
actually the majority of the algorithm is just doing things from the FA and determining error via the swap method described in the paper.
I don't think I will have time to talk about the final algorithm for learning the parameter vector $\bv{\lambda}$. So I think my presentation will look like
\begin{enumerate}
  \item A flowchart, maybe across multiple slides, describing the FA algorithm (maybe retermed because it doesn't seem to really be fitting).
        In these slides I will define terms that will be unfamiliar to most people and try to explain the machinery that this algorithm develops, with a focus
        on the machinery that is used in the EEA.
  \item Explain the first part of the EEA, which is usage of the machinery developed in the FA.
  \item Explain the swap-test based error estimation. A figure (maybe animated to show sampling progression and illustrate the approximate number of required steps)
        will be much more useful here than me trying to explain the process verbally I think.
  \item Discuss the limitations section. If I have the chance to read the later papers on the topic that cite this one I will try to mention that some of
        these are no longer limitations i.e. \url{https://doi.org/10.1016/j.physa.2023.128521} removes the need for a sparse matrix, greatly broadening applications.
        Maybe talk about sparse matrices in more detail and why it is very useful to be able to go beyond them (I don't really know this yet)
\end{enumerate}
I feel that this is quite a lot of content and I will learn what I have to cut down on when I start on my slides. I do think though that I will manage
to get across sufficient information in the given time that people will at least have an idea of what is going on.
\end{document}